\scanpage{27}

\begin{center}-- 216 --\end{center}
\hfill n$^{\circ}$ 391

(\S\ 8, n$^{\circ}$ 3, formule (2)) ; autrement dit, la suite des exposants est :
\[
1,\ 5,\ 7,\ 9,\ 11,\ 13,\ 17.
\]

(X)\quad On déduit de (IX) et du \S\ 8, n$^{\circ}$ 3, cor. 1 de la prop. 3, que l'ordre du
groupe de Weyl $W(R)$ est :
\[
2.6.8.10.12.14.18=2^{10}.3^4.5.7.
\]

(XI) Le même raisonnement qu'au n$^{\circ}$ 8 montre que le seul automorphisme du schéma
de Dynkin $E_7$ est l'identité.

12. \uline{Système de type} $E_6$.

(I)\quad Gardons toujours les notations du n$^{\circ}$ 10, et considérons l'intersection $R''$
de $R$ et du sous-espace $V''$ de codimension 2, orthogonal aux deux vecteurs $\varepsilon_8+\varepsilon_7$
et $\varepsilon_7-\varepsilon_6$ ; ce sont donc les vecteurs $x=\sum_{i=1}^8\xi_i\varepsilon_i$ de $R$ tels que
$\xi_8=-\xi_7=-\xi_6$;
autrement dit :
\[
\pm\varepsilon_i\pm\varepsilon_j,\quad(i\ne j,\ 1\leq i\leq5,\ 1\leq j\leq5)
\]
\[
\pm\frac12\left(\varepsilon_8-\varepsilon_7-\varepsilon_6+
\sum_{i=1}^5(-1)^{\unsure{\nu|\upsilon}(i)}\varepsilon_i\right)
\quad\text{où la somme }\sum_{i=1}^5\unsure{\nu|\upsilon}(i)\text{ est}
\]
\uline{paire}.

Il est clair que $R''$ engendre $V''$, donc est un système de racines réduit de
rang 6, ayant $n=\binom52.4+2^5=72$ éléments.

(II)\quad Si l'on exprime qu'une racine $\sum_{i=1}^8n_i\alpha_i$ de $R$ est orthogonale à la fois à
$\varepsilon_8+\varepsilon_7$ et à $\varepsilon_7-\varepsilon_6$, il vient, en vertu du n$^{\circ}$ 10, (II), $n_8=n_7=0$. Le même
raisonnement que dans le n$^{\circ}$ 11 montre alors que les $\alpha_i$ pour $1\leq i\leq6$ forment
un système de \uline{racines simples} $S''$ de $R''$. Le calcul des $(\alpha_i|\alpha_j)$ fait au n$^{\circ}$ 10
montre que le schéma de Dynkin de $R''$ est de type $E_6$, et $R''$ est irréductible.

(III) En vertu du \S\ 8, n$^{\circ}$ 3, formule (4), on a $h=n/\ell=12$.

\mspage{26}
